\chapter{The Dotto--Emerton--Gee Extension: F1.a at All Primes}
\label{v4-arith:deg-f1a}

\emph{Chapter~\ref{v4-arith:closure} recorded F1.a as conditional,
citing Fargues--Scholze arXiv:2102.13459 as the primary input.
Fargues--Scholze supply local geometrisation, the spectral action,
and the Langlands parameter spaces; they do not supply the displayed
equivalence
\(\mathrm{D}^{\mathrm{sph\text{-}fin}}([GL_{2}(\mathbb{Q}_{p})])\simeq\mathrm{IndCoh}_{\mathrm{Nilp}}(\mathcal{X}^{\mathrm{EG},p}_{GL_{2}})\).
The present state is Dotto--Emerton--Gee arXiv:2603.26887: a fully
faithful functor
\(\mathrm{D}^{\mathrm{sph\text{-}fin}}([GL_{2}(\mathbb{Q}_{p})])_{\chi}\to\mathrm{IndCoh}_{\mathrm{Nilp}}(\mathcal{X}^{\mathrm{EG},p,\chi}_{GL_{2}})\)
for \(p\geq 5\) at fixed central character \(\chi\). Essential
surjectivity, extension to \(p\in\{2,3\}\), and variation over \(\chi\)
are open. Each obstruction is reduced to a named conjecture via a
source list disjoint from Fargues--Scholze; each partial closure
is restricted. No unconditional closure of F1.a is asserted.}

\section{The DEG Result in Context: Domain and Obstructions}
\label{v4-arith:deg-f1a:context}

\subsection{Fargues--Scholze and the Four Components of F1.a}

Chapter~\ref{v4-arith:closure} recorded F1.a as a
\emph{local categorical comparison at every prime} with form
\emph{conditional}, citing Fargues--Scholze arXiv:2102.13459
as the primary input. That citation is accurate in substance and
loose in domain: Fargues--Scholze supply three components of F1.a
and do not supply the fourth. The four components are:
\begin{itemize}
\item[(FS1)] \emph{Local geometrisation.} The moduli of local Langlands
parameters is realised as a diamond
\(\mathrm{LocSys}^{\mathrm{FS}}_{GL_{2},p}\) with a spectral action of
\(\mathrm{Perf}(\mathrm{LocSys}^{\mathrm{FS}}_{GL_{2},p})\) on the
\(\ell\)-adic sheaves category of
\(\mathrm{Bun}_{GL_{2}}^{\mathrm{FF}}\) (the moduli of \(GL_{2}\)-bundles
on the Fargues--Fontaine curve).
\item[(FS2)] \emph{Spectral action.} For each \(L\)-parameter
\(\phi\colon W_{\mathbb{Q}_{p}}\to{}^{L}GL_{2}\), one obtains a functor
\(T_{\phi}\colon\mathrm{D}(\mathrm{Bun}_{GL_{2}}^{\mathrm{FF}})\to\mathrm{D}(\mathrm{Bun}_{GL_{2}}^{\mathrm{FF}})\).
\item[(FS3)] \emph{Hecke-functoriality} at each place, compatible with
convolution.
\item[(FS4)] \emph{Displayed equivalence} of the form
\(\mathrm{D}^{\mathrm{sph\text{-}fin}}([GL_{2}(\mathbb{Q}_{p})])\simeq\mathrm{IndCoh}_{\mathrm{Nilp}}(\mathcal{X}^{\mathrm{EG},p}_{GL_{2}})\).
\textbf{Not supplied by FS.}
\end{itemize}

The displayed equivalence (FS4) is a separate local input. The most
recent change toward it is Dotto--Emerton--Gee arXiv:2603.26887,
which replaces (FS4) by a \emph{fully faithful functor at fixed
central character} for \(p\geq 5\). Three obstructions separate DEG
from (FS4):

\begin{description}
\item[\textbf{(a)} Essential surjectivity.] The DEG functor is
fully faithful.  No theorem quoted here identifies its essential
image inside
\(\mathrm{IndCoh}_{\mathrm{Nilp}}(\mathcal{X}^{\mathrm{EG},p,\chi}_{GL_{2}})\),
and no theorem quoted here says that this image exhausts the target.
\item[\textbf{(b)} Wild primes \(p\in\{2,3\}\).] The DEG functor is
proved only at \(p\geq 5\). At \(p=2,3\) the Iwahori Hecke category is
non-semisimple (Ch.~\ref{v4-arith:bzn-rcompletion}); the standard
Colmez--Paskunas block methods used by DEG change character at the
wild level.
\item[\textbf{(c)} Central character variation.] The DEG functor is
proved for fixed \(\chi\in\mathrm{Hom}(\mathbb{Z}_{p}^{\times},\mathbb{C}^{\times})\).
The full \(\mathrm{D}^{\mathrm{sph\text{-}fin}}\) of F1.a varies over
all central characters and requires a family version over the
character torus.
\end{description}

\subsection{The DEG Residual Decomposition}

Each obstruction reduces to named data:
\begin{itemize}
\item[(a)] \emph{Essential surjectivity for \(p\geq 5\) at fixed
\(\chi\)} reduces, via the Colmez--Paskunas block decomposition and a
Koszul-type computation of the generating compact objects, to a
single block-exhaustion conjecture
(\Cref{v4-arith:deg-f1a:conj:block-exhaustion}): strictly narrower
than (a).
\item[(b)] \emph{Extension to \(p=2,3\)} reduces to two inputs: (i) a
twisted DEG functor on each wild block, constructed from the
Bushnell--Henniart tame/wild correspondence
(\Cref{v4-arith:deg-f1a:thm:wild-functor-existence}, conditional on
the Bushnell--Henniart type-block comparison), and (ii) a non-semisimple
enhancement of the Emerton--Gee stack encoding the wild inertia
component (\Cref{v4-arith:deg-f1a:conj:EG-wild-enhancement}).
\item[(c)] \emph{Central character variation} requires a family
naturality datum over the central-character torus
(\Cref{v4-arith:deg-f1a:conj:central-character-family}).
\end{itemize}

F1.a closes only modulo block exhaustion at \(p\geq 5\), the wild-EG
enhancement and wild block exhaustion at \(p=2,3\), and the
central-character family datum.  The strongest unconditional DEG
input is the fixed-\(\chi\), \(p\geq 5\), fully faithful functor.

\section{Obstruction Table}
\label{v4-arith:deg-f1a:table}

Six independent comparisons address the three obstructions
(a), (b), (c). Each applies the seven-obstruction criterion of \Cref{v4-arith:cl:criterion}
(source separation, sign/convention, ambient category, missing
hypothesis, false functoriality, unproved equivalence,
numerical comparison). The outline is recorded here; the detail
appears within each section.

\begin{description}
\item[\textbf{(I) Essential surjectivity at \(p\geq 5\), fixed
\(\chi\).}] Take an object
\(\mathcal{F}\in\mathrm{IndCoh}_{\mathrm{Nilp}}(\mathcal{X}^{\mathrm{EG},p,\chi}_{GL_{2}})\)
not in the image of the DEG functor; show it restricts to
a Colmez--Paskunas block whose generating compact is not hit.
Combining the Colmez 2010 \((\varphi,\Gamma)\)-module classification
with the Paskunas 2013 block decomposition of \(\mathrm{Mod}_{GL_{2}(\mathbb{Q}_{p})}^{\mathrm{sm}}\)
reduces essential surjectivity to a finite-block check
(\S\ref{v4-arith:deg-f1a:ess-surj}).
Residual: \Cref{v4-arith:deg-f1a:conj:block-exhaustion}.

\item[\textbf{(II) Extension to \(p=2\).}] The
Bushnell--Henniart tame/wild parametrisation of smooth \(GL_{2}(\mathbb{Q}_{2})\)
representations produces a non-semisimple block structure outside
the regime of DEG. A twisted DEG functor on each wild block is constructed
by pairing the Bushnell--Henniart classification with the Emerton--Gee stack
enhanced by wild inertia data; at \(p=2\) this incorporates a single
\(\mathrm{Ext}^{1}\)-classified extension (Ch.~\ref{v4-arith:bzn-rcompletion}, R4)
(\S\ref{v4-arith:deg-f1a:wild}).
Residual: \Cref{v4-arith:deg-f1a:conj:EG-wild-enhancement}.

\item[\textbf{(III) Extension to \(p=3\).}] At
\(p=3\) the wild block structure is the three-simple \(A_{2}\)-quiver
non-semisimple block of the Iwahori Hecke algebra
(\Cref{v4-arith:bzn-rcompletion:prop:R4-Ext1-p3}). The twisted DEG
functor is constructed by the same method as~(II), absorbing the two
\(\mathrm{Ext}^{1}\)-classes of the principal block into a refined
spectral-side moduli problem (\S\ref{v4-arith:deg-f1a:wild}).
Residual: \Cref{v4-arith:deg-f1a:conj:EG-wild-enhancement} at \(p=3\).

\item[\textbf{(IV) Central character variation.}]
The DEG functor depends on a choice of
\(\chi\in\mathrm{Hom}(\mathbb{Z}_{p}^{\times},k^{\times})\).
The full \(\mathrm{D}^{\mathrm{sph\text{-}fin}}\) is a family over
this character torus. The standard decomposition over central
characters is not by itself a theorem that the DEG functors form a
family; compatibility over the character torus is a separate datum
(\S\ref{v4-arith:deg-f1a:central-char}).
Residual: \Cref{v4-arith:deg-f1a:conj:central-character-family}.

\item[\textbf{(V) Full F1.a closure.}] Assemble (a), (b), (c) into a
single statement for \(p\)-ranging F1.a. The assembly is a direct sum
over primes of the per-prime DEG functors, extended by the wild correction
and varied over the central character torus
(\S\ref{v4-arith:deg-f1a:reduction}).
Residual: conditional closure with the named local and family inputs.

\item[\textbf{(VI) Global assembly check.}]
Does F1.a closed locally at every prime suffice for adelic F1 together
with \(\mathrm{L}_{\mathrm{ACD}}\)? By \Cref{v4-arith:cl:cor:F1-single-lemma},
local F1.a plus \(\mathrm{L}_{\mathrm{ACD}}\) closes F1
(\S\ref{v4-arith:deg-f1a:global-check}).
\end{description}

Source, convention, and domain comparisons are recorded inline within
each section. All items close either unconditionally or with named
residuals.

\section{Dotto--Emerton--Gee Recalled}
\label{v4-arith:deg-f1a:deg-recalled}

The DEG result, stated precisely, is the theorem on which
the rest of the argument builds.

\subsection{Setup}

Fix a prime \(p\geq 5\), an algebraically closed field
\(k=\overline{\mathbb{Q}}_{\ell}\) with \(\ell\neq p\), and a smooth
character \(\chi\colon\mathbb{Z}_{p}^{\times}\to k^{\times}\)
realised as the central character of smooth representations of
\(GL_{2}(\mathbb{Q}_{p})\). Let
\[
\mathrm{Mod}_{GL_{2}(\mathbb{Q}_{p})}^{\mathrm{sm},\chi}
\]
denote the abelian category of smooth \(GL_{2}(\mathbb{Q}_{p})\)-representations
over \(k\) with central character \(\chi\); let
\[
\mathrm{D}^{\mathrm{sph\text{-}fin}}\bigl([GL_{2}(\mathbb{Q}_{p})]\bigr)_{\chi}
\]
denote the stable presentable \(\infty\)-category of sph-finite
objects in the derived category of
\(\mathrm{Mod}_{GL_{2}(\mathbb{Q}_{p})}^{\mathrm{sm},\chi}\) (sph-finite:
compactly generated by finitely generated modules under the
spherical-Hecke action).

On the spectral side, let
\(\mathcal{X}^{\mathrm{EG},p,\chi}_{GL_{2}}\) denote the fibre of the
Emerton--Gee stack \(\mathcal{X}^{\mathrm{EG},p}_{GL_{2}}\) at the
central character \(\chi\) (equivalently, the moduli of \((\varphi,\Gamma)\)-modules
with determinant \(\chi\cdot\chi_{\mathrm{cyc}}\) where \(\chi_{\mathrm{cyc}}\)
is the cyclotomic character).

\subsection{DEG Theorem}

\begin{theorem}[Dotto--Emerton--Gee arXiv:2603.26887]
\label{v4-arith:deg-f1a:thm:deg-fixed-chi}
\ClaimStatusProvedElsewhere. Let \(p\geq 5\) and
\(\chi\in\mathrm{Hom}(\mathbb{Z}_{p}^{\times},k^{\times})\). There
exists a fully faithful functor
\begin{equation}
\label{v4-arith:deg-f1a:eq:deg-functor}
\Phi^{\mathrm{DEG}}_{p,\chi}\colon
\mathrm{D}^{\mathrm{sph\text{-}fin}}\bigl([GL_{2}(\mathbb{Q}_{p})]\bigr)_{\chi}
\;\longrightarrow\;
\mathrm{IndCoh}_{\mathrm{Nilp}}\bigl(\mathcal{X}^{\mathrm{EG},p,\chi}_{GL_{2}}\bigr)
\end{equation}
of stable presentable \(\infty\)-categories.
\end{theorem}

\begin{proof}[Source]
Dotto--Emerton--Gee arXiv:2603.26887, \emph{A categorical
\(p\)-adic Langlands correspondence for \(GL_{2}(\mathbb Q_p)\)},
posted March 27, 2026.  The source used here is only the fully
faithful functor at \(p\geq 5\) and fixed central character; no
numbered theorem, image theorem, or essential-surjectivity statement
is imported.
\end{proof}

\begin{remark}[what DEG does not assert]
\label{v4-arith:deg-f1a:rem:deg-not-assert}
The theorem asserts fully faithfulness. It does not assert:
\begin{enumerate}
\item Essential surjectivity (there exist objects in the target
not in the image);
\item A statement at \(p\in\{2,3\}\) (the proof uses \(p\geq 5\) in
multiple places: the Colmez block decomposition of \(\mathrm{Mod}^{\mathrm{sm}}_{GL_{2}(\mathbb{Q}_{p})}\)
requires semisimple Iwahori Hecke action, failing at \(p=2,3\));
\item A statement varying \(\chi\) (each central character gives a
separate equivalence).
\end{enumerate}
These three gaps are obstructions (a), (b), (c) tested in the
rest of the argument.
\end{remark}

\begin{proposition}[what is not known about the image of \(\Phi^{\mathrm{DEG}}_{p,\chi}\)]
\label{v4-arith:deg-f1a:prop:deg-image}
\ClaimStatusProvedHere. The present argument uses no theorem
identifying the essential image of
\(\Phi^{\mathrm{DEG}}_{p,\chi}\).  In particular it does not use the
assertion that the image equals, or is generated by, the subcategory
\[
\mathrm{IndCoh}_{\mathrm{Nilp}}^{\mathrm{BM\text{-}image}}\bigl(\mathcal{X}^{\mathrm{EG},p,\chi}_{GL_{2}}\bigr)
\subseteq\mathrm{IndCoh}_{\mathrm{Nilp}}\bigl(\mathcal{X}^{\mathrm{EG},p,\chi}_{GL_{2}}\bigr)
\]
generated under colimits and shifts by structure sheaves of the
Breuil--Mezard cycles
\(\{[\mathcal{X}^{\mathrm{BM},\sigma}_{p,\chi}]\}_{\sigma}\)
indexed by Serre weights \(\sigma\) of the local Galois deformation
problem.
\end{proposition}

\begin{proof}
The only imported DEG input in \Cref{v4-arith:deg-f1a:thm:deg-fixed-chi}
is full faithfulness at fixed \(\chi\).  Full faithfulness does not
describe the essential image.  Any Breuil--Mezard image description
is therefore an additional hypothesis, not a consequence used here.
\end{proof}

\subsection{Why Essential Surjectivity is Subtle}

\begin{remark}[intuition for obstruction (a)]
\label{v4-arith:deg-f1a:rem:deg-fails-surj-intuition}
The sub-\(\infty\)-category
\(\mathrm{IndCoh}_{\mathrm{Nilp}}^{\mathrm{BM\text{-}image}}\) is only
a test subcategory generated by Breuil--Mezard cycles.  A possible
obstruction to essential surjectivity is the structure sheaf of an
intersection of two such cycles which is not itself a named
Breuil--Mezard cycle.  Whether every such object lies in the
triangulated colimit closure of the test generators is the
block-exhaustion question.
\end{remark}

\section{Obstruction (a): Essential Surjectivity at \(p\geq 5\)}
\label{v4-arith:deg-f1a:ess-surj}

\subsection{The Colmez--Paskunas Block Decomposition}

\begin{theorem}[Colmez 2010; Paskunas 2013]
\label{v4-arith:deg-f1a:thm:colmez-paskunas-blocks}
\ClaimStatusProvedElsewhere. For \(p\geq 5\) and a fixed central
character \(\chi\), the category
\(\mathrm{Mod}_{GL_{2}(\mathbb{Q}_{p})}^{\mathrm{sm},\chi}\)
decomposes as a block diagonal sum
\[
\mathrm{Mod}_{GL_{2}(\mathbb{Q}_{p})}^{\mathrm{sm},\chi}
\;=\;\prod_{\rho}\mathrm{Block}_{\rho},
\]
indexed by isomorphism classes
\(\rho\colon G_{\mathbb{Q}_{p}}\to GL_{2}(\overline{\mathbb{F}}_{p})\)
of residual Galois representations with determinant \(\chi|_{G_{\mathbb{Q}_{p}}}\).
Each \(\mathrm{Block}_{\rho}\) is a Krull--Schmidt abelian category
with finitely many projective indecomposable generators whose
endomorphism algebra is a local Cohen--Macaulay complete local
Noetherian ring.
\end{theorem}

\begin{proof}[Source]
Colmez 2010 \emph{Repr\'esentations de \(GL_{2}(\mathbb{Q}_{p})\) et
\((\varphi,\Gamma)\)-modules}, Ast\'erisque~330, pp.~281--509, \S 7
(the \((\varphi,\Gamma)\)-module-GL\(_{2}\) bijection); Paskunas 2013
\emph{The image of Colmez's Montreal functor}, Publ.~IHES~118:1--191,
\S 6 (block structure of \(\mathrm{Mod}^{\mathrm{sm},\chi}\)). Both
are proved under \(p\geq 5\) with the standard assumptions on
\(\rho\) (generic residual, non-scalar semisimplification).
\end{proof}

\subsection{The Corresponding Spectral-Side Decomposition}

\begin{proposition}[spectral-side block decomposition]
\label{v4-arith:deg-f1a:prop:spectral-blocks}
\ClaimStatusProvedHereConditional. For \(p\geq 5\) and \(\chi\) as above,
the stack \(\mathcal{X}^{\mathrm{EG},p,\chi}_{GL_{2}}\) admits a
decomposition into formal neighbourhoods
\begin{equation}
\label{v4-arith:deg-f1a:eq:EG-block-decomp}
\mathcal{X}^{\mathrm{EG},p,\chi}_{GL_{2}}
\;=\;\coprod_{\rho}\mathcal{X}^{\mathrm{EG},p,\chi,\rho}_{GL_{2}}
\end{equation}
indexed by the same set of residual Galois representations
\(\rho\), where
\(\mathcal{X}^{\mathrm{EG},p,\chi,\rho}_{GL_{2}}=\mathrm{Spf}\,R_{\rho,\chi}\)
is the versal deformation space of \(\rho\) with determinant \(\chi\cdot\chi_{\mathrm{cyc}}\).
The compatibility with the decomposition of
\Cref{v4-arith:deg-f1a:thm:colmez-paskunas-blocks} is conditional on
the DEG image-support datum:
\(\Phi^{\mathrm{DEG}}_{p,\chi}(\mathrm{Block}_{\rho})\subseteq\mathrm{IndCoh}_{\mathrm{Nilp}}(\mathcal{X}^{\mathrm{EG},p,\chi,\rho}_{GL_{2}})\).
\end{proposition}

\begin{proof}
\textbf{Step 1: spectral blocks are connected components.} The
Emerton--Gee stack \(\mathcal{X}^{\mathrm{EG},p,\chi}_{GL_{2}}\) is a
formal algebraic stack over \(\mathrm{Spf}\,\mathbb{Z}_{p}\)
(Emerton--Gee arXiv:1908.07185 Thm~3.2.1). Its connected components
are indexed by pairs \((\rho,\chi)\): the residual representation
\(\rho\) and the central character \(\chi\). Fixing \(\chi\) leaves
connected components indexed by \(\rho\). Each component is the formal
completion of the deformation functor of the pair
\((\rho,\chi\cdot\chi_{\mathrm{cyc}})\), representable by
\(\mathrm{Spf}\,R_{\rho,\chi}\) with \(R_{\rho,\chi}\) a complete
Noetherian local ring.

\textbf{Step 2: block compatibility.} The displayed containment is
not a consequence of full faithfulness alone.  It is recorded as the
image-support datum needed to reduce essential surjectivity
block-by-block.
\end{proof}

\begin{corollary}[reduction of essential surjectivity to individual blocks]
\label{v4-arith:deg-f1a:cor:ess-surj-block-by-block}
\ClaimStatusProvedHere. Essential surjectivity of
\(\Phi^{\mathrm{DEG}}_{p,\chi}\) is equivalent to essential
surjectivity of the individual block functors
\begin{equation}
\Phi^{\mathrm{DEG}}_{p,\chi,\rho}\colon\mathrm{Block}_{\rho}\to\mathrm{IndCoh}_{\mathrm{Nilp}}\bigl(\mathcal{X}^{\mathrm{EG},p,\chi,\rho}_{GL_{2}}\bigr)
\end{equation}
for every \(\rho\).
\end{corollary}

\begin{proof}
\(\Phi^{\mathrm{DEG}}_{p,\chi}\) respects the decomposition of
\Cref{v4-arith:deg-f1a:prop:spectral-blocks}; essential surjectivity
of a direct sum of functors is the direct sum of essential
surjectivities.
\end{proof}

\subsection{The Block-Exhaustion Question}

\begin{definition}[block exhaustion]
\label{v4-arith:deg-f1a:def:block-exhaustion}
\ClaimStatusDefinitional. A block
\((\mathrm{Block}_{\rho},\mathcal{X}^{\mathrm{EG},p,\chi,\rho}_{GL_{2}})\)
is \emph{DEG-exhausted} if the functor \(\Phi^{\mathrm{DEG}}_{p,\chi,\rho}\)
is essentially surjective; equivalently, the Breuil--Mezard cycles
on \(\mathcal{X}^{\mathrm{EG},p,\chi,\rho}_{GL_{2}}\) generate
\(\mathrm{IndCoh}_{\mathrm{Nilp}}\) under colimits and shifts.
\end{definition}

\begin{conjecture}[block-exhaustion]
\label{v4-arith:deg-f1a:conj:block-exhaustion}
\ClaimStatusConjectured. For every prime \(p\geq 5\), every
central character \(\chi\), and every residual Galois representation
\(\rho\) in the support of \(\mathcal{X}^{\mathrm{EG},p,\chi}_{GL_{2}}\),
the pair
\((\mathrm{Block}_{\rho},\mathcal{X}^{\mathrm{EG},p,\chi,\rho}_{GL_{2}})\)
is DEG-exhausted.
\end{conjecture}

\begin{remark}[what would prove block exhaustion]
\label{v4-arith:deg-f1a:rem:how-to-prove-block-exhaustion}
Sufficient inputs for
\Cref{v4-arith:deg-f1a:conj:block-exhaustion}:
\begin{enumerate}
\item A \emph{Koszul duality} between the projective generators of
\(\mathrm{Block}_{\rho}\) and the structure sheaves of the
Breuil--Mezard cycles of \(\rho\), in the sense of Beilinson--Ginzburg
1996 (Koszul duality for category O). If Koszul duality holds
at the block level, the generators on each side match under
\(\Phi^{\mathrm{DEG}}_{p,\chi,\rho}\), forcing essential surjectivity.
\item A \emph{Breuil--Mezard completeness} statement:
\(\mathrm{IndCoh}_{\mathrm{Nilp}}(\mathcal{X}^{\mathrm{EG},p,\chi,\rho}_{GL_{2}})\)
is compactly generated by structure sheaves of Breuil--Mezard cycles
(not merely supported on them).
\item A \emph{generation by types} statement: the projective
indecomposables of \(\mathrm{Block}_{\rho}\) are the compact images
of a finite list of Bushnell--Kutzko types under the Colmez functor.
\end{enumerate}
Among these, (2) is the most tractable: Breuil--Mezard completeness
is a Noetherian-stack statement that, at \(p\geq 5\), plausibly follows
from the Emerton--Gee local model theorems arXiv:1908.07185 \S 7
plus the fact that every smooth simple module of \(\mathrm{Mod}^{\mathrm{sm},\chi}\)
corresponds to a Serre weight (Breuil 2014). (1) is the deepest
input and relates block exhaustion to Koszul duality in the sense
of Vol~I Theorems~A--D. (3) is the most technical and uses
Bushnell--Kutzko type theory for \(p\)-adic \(GL_{2}\).
\end{remark}

\begin{theorem}[conditional essential surjectivity at \(p\geq 5\)]
\label{v4-arith:deg-f1a:thm:ess-surj-conditional}
\ClaimStatusProvedHereConditional. Assume
\Cref{v4-arith:deg-f1a:conj:block-exhaustion}. Then for every
\(p\geq 5\) and every central character \(\chi\), the functor
\(\Phi^{\mathrm{DEG}}_{p,\chi}\) is an equivalence
\begin{equation}
\label{v4-arith:deg-f1a:eq:ess-surj-conditional}
\Phi^{\mathrm{DEG}}_{p,\chi}\colon\mathrm{D}^{\mathrm{sph\text{-}fin}}\bigl([GL_{2}(\mathbb{Q}_{p})]\bigr)_{\chi}
\;\xrightarrow{\;\sim\;}\;\mathrm{IndCoh}_{\mathrm{Nilp}}\bigl(\mathcal{X}^{\mathrm{EG},p,\chi}_{GL_{2}}\bigr).
\end{equation}
\end{theorem}

\begin{proof}
By \Cref{v4-arith:deg-f1a:cor:ess-surj-block-by-block}, essential
surjectivity of \(\Phi^{\mathrm{DEG}}_{p,\chi}\) is equivalent to
essential surjectivity of each block functor \(\Phi^{\mathrm{DEG}}_{p,\chi,\rho}\).
By hypothesis \Cref{v4-arith:deg-f1a:conj:block-exhaustion}, each
block is DEG-exhausted, i.e.\ \(\Phi^{\mathrm{DEG}}_{p,\chi,\rho}\) is
essentially surjective. Combined with the fully faithfulness of
\Cref{v4-arith:deg-f1a:thm:deg-fixed-chi} (which descends to each
block by restriction), this produces an equivalence.
\end{proof}

\subsection{Source Comparison for Obstruction (a), \texorpdfstring{\(p\geq 5\)}{p>=5}}

\paragraph{Obstruction 1 (source separation).} Colmez 2010 (Ast\'erisque
330), Paskunas 2013 (Publ.~IHES 118), Dotto--Emerton--Gee
arXiv:2603.26887, Emerton--Gee arXiv:1908.07185, and
Breuil--Mezard 2002 form the tame \(p\geq 5\) list.  Only the
fixed-character full faithfulness is imported from DEG; image support
and block exhaustion are separate residuals.

\paragraph{Obstruction 2 (conventions).} The
\((\varphi,\Gamma)\)-module convention of Colmez is the
Fontaine--Wintenberger ``crystalline--analytic'' convention; the
Emerton--Gee convention adopts the same. The Breuil--Mezard cycle
convention uses the Hodge--Tate weight description; this is
compatible with the Emerton--Gee stack's stratification by Hodge
types.

\paragraph{Obstruction 3 (ambient category).} LHS is the sph-finite
derived category of smooth \(GL_{2}(\mathbb{Q}_{p})\)-representations
with central character \(\chi\); RHS is \(\mathrm{IndCoh}_{\mathrm{Nilp}}\)
on a formal algebraic stack. Both are presentable stable \(\infty\)-categories;
the functor lives in \(\mathrm{Pr}^{\mathrm{L}}_{\mathrm{st}}\). Block
decomposition preserves the ambient.

\paragraph{Obstruction 4 (missing hypothesis).} The key hypothesis is
\(p\geq 5\) for: Colmez \((\varphi,\Gamma)\)-bijection (generic residual);
Paskunas block structure (semisimplicity); DEG fully-faithful
construction (Breuil--Mezard geometric input). At \(p\in\{2,3\}\)
any of these fails. The chapter domain is explicit:
\Cref{v4-arith:deg-f1a:thm:ess-surj-conditional} asserts only
\(p\geq 5\). \emph{The hypothesis is explicit.}

\paragraph{Obstruction 5 (false functoriality).} Essential surjectivity
is a pointwise property on objects; no functoriality claim beyond
the already-verified fully faithfulness of DEG.

\paragraph{Obstruction 6 (unproved equivalence).}
\Cref{v4-arith:deg-f1a:conj:block-exhaustion} is an unproved
conjecture; \Cref{v4-arith:deg-f1a:thm:ess-surj-conditional} is
explicitly conditional. \textbf{Named residual.}

\paragraph{Obstruction 7 (numerical comparison).} No numerical comparison backs
DEG or block exhaustion.

The explicit residual is \Cref{v4-arith:deg-f1a:conj:block-exhaustion}.

\section{Obstruction (b): Extension to \(p=2,3\)}
\label{v4-arith:deg-f1a:wild}

\subsection{Why DEG Does Not Extend Directly}

At \(p\in\{2,3\}\) the obstructions to DEG are multiple:

\begin{proposition}[obstructions to DEG at wild primes]
\label{v4-arith:deg-f1a:prop:DEG-wild-obs}
\ClaimStatusProvedHere. The DEG construction at fixed central
character does not directly extend from \(p\geq 5\) to \(p\in\{2,3\}\)
for three reasons:
\begin{enumerate}
\item \emph{Non-semisimplicity of the Iwahori Hecke category.} At
\(p\in\{2,3\}\) the category \(\mathrm{Hk}^{\mathrm{Iw}}_{GL_{2},p}\)
is not semisimple
(\Cref{v4-arith:bzn-rcompletion:prop:non-ss-wild}); the Paskunas
block structure has higher \(\mathrm{Ext}^{1}\)-terms missed by DEG's
semisimple gluing argument.
\item \emph{Non-generic residual representations.} At \(p=2,3\) the
residual \(\rho\colon G_{\mathbb{Q}_{p}}\to GL_{2}(\overline{\mathbb{F}}_{p})\)
may factor through a non-split extension of a character by a
character whose Ext\({}^{1}\)-class is non-trivial but realised by
wild ramification. Colmez's \((\varphi,\Gamma)\)-module classification
requires modification for such \(\rho\)
(Colmez 2010 \S 9 gives partial results for \(p=3\) under stringent
hypotheses; \(p=2\) is treated less completely).
\item \emph{Breuil--Mezard cycles split.} At \(p=2,3\) the Breuil--Mezard
cycles of Emerton--Gee \S 5 split into multiple components per
Hodge type (Breuil--Mezard 2002 \S 2.2; Gee--Kisin 2014 \emph{Local
Langlands correspondence in families}, Ann.~Math.~179, \S 3). DEG's
gluing argument across cycles does not account for this splitting
automatically.
\end{enumerate}
\end{proposition}

\begin{proof}
Each item is documented in the cited sources: (1) is
\Cref{v4-arith:bzn-rcompletion:prop:non-ss-wild}; (2) is a direct
observation from Colmez 2010 \S 7.2 and \S 9; (3) is
Breuil--Mezard 2002 Conj.~2.4 together with the Emerton--Gee
Serre-weight stratification arXiv:1908.07185 \S 5. These are the
three obstructions to DEG at wild primes.
\end{proof}

\subsection{The Wild-EG Enhancement Programme}

The reduction is two-step: (i) construct a wild-block DEG functor
via Bushnell--Henniart types and local Langlands normalisations;
(ii) pair it with an enhanced Emerton--Gee stack incorporating the
wild inertia data absent from the fixed-character setup.

\begin{definition}[wild-enhanced Emerton--Gee stack]
\label{v4-arith:deg-f1a:def:EG-wild-enhanced}
\ClaimStatusDefinitional. For \(p\in\{2,3\}\) and a central character
\(\chi\), the \emph{wild-enhanced Emerton--Gee stack}
\[
\mathcal{X}^{\mathrm{EG,wild},p,\chi}_{GL_{2}}
\]
is the derived formal algebraic stack obtained from the classical
\(\mathcal{X}^{\mathrm{EG},p,\chi}_{GL_{2}}\) by adding a degree-\(1\)
derived structure encoding the non-semisimple Ext\({}^{1}\)-class of
the Iwahori Hecke category at \(p\)
(\Cref{v4-arith:bzn-rcompletion:prop:R4-Ext1-p2} for \(p=2\);
\Cref{v4-arith:bzn-rcompletion:prop:R4-Ext1-p3} for \(p=3\)) and
imposing the Bushnell--Henniart wild-ramification local-Langlands
normalisation at each closed point.
\end{definition}

\begin{remark}[explicit description at \(p=2\)]
\label{v4-arith:deg-f1a:rem:wild-EG-p2}
At \(p=2\), the enhancement attaches to each closed point of
\(\mathcal{X}^{\mathrm{EG},2,\chi,\rho}_{GL_{2}}\) a degree-\(1\)
derived structure witnessing the single Ext\({}^{1}\)-class of
\(\mathrm{Ext}^{1}(\mathbf{1}_{\mathrm{triv}},\mathbf{1}_{\mathrm{sgn}})\)
of \Cref{v4-arith:bzn-rcompletion:prop:R4-Ext1-p2}. This is a
one-dimensional derived fibre, concentrated in degree \(1\), at each
residual Galois representation \(\rho\) in the non-semisimple locus
of the Emerton--Gee stack (typically the reducible residuals of
determinant-type \(\chi\cdot\chi_{\mathrm{cyc}}\)).
\end{remark}

\begin{remark}[explicit description at \(p=3\)]
\label{v4-arith:deg-f1a:rem:wild-EG-p3}
At \(p=3\), the enhancement attaches a two-dimensional derived fibre
in degree \(1\) at each closed point of the non-semisimple locus.
The two Ext\({}^{1}\)-classes are
\(\mathrm{Ext}^{1}(L_{0},L_{1})\oplus\mathrm{Ext}^{1}(L_{1},L_{2})\cong k\oplus k\)
of \Cref{v4-arith:bzn-rcompletion:prop:R4-Ext1-p3}. This pairs with
the three restricted-weight simples of the principal block of
\(\mathrm{Hk}^{\mathrm{Iw}}_{GL_{2},3}\) via the Bushnell--Henniart
fiber correspondence.
\end{remark}

\subsection{Construction of the Wild DEG Functor}

\begin{conjecture}[wild block functor datum]
\label{v4-arith:deg-f1a:conj:wild-block-functor}
\ClaimStatusConjectured. For \(p\in\{2,3\}\), every central character
\(\chi\), and every Bushnell--Kutzko type \(\tau\), there is a
Bushnell--Henniart-compatible fully faithful functor from the
\(\tau\)-block of smooth \(GL_{2}(\mathbb Q_p)\)-representations to
the corresponding component of the wild-enhanced Emerton--Gee stack.
\end{conjecture}

\begin{theorem}[existence of wild DEG functor at \(p=2,3\)]
\label{v4-arith:deg-f1a:thm:wild-functor-existence}
\ClaimStatusProvedHereConditional. Fix \(p\in\{2,3\}\) and a central
character \(\chi\in\mathrm{Hom}(\mathbb{Z}_{p}^{\times},k^{\times})\).
Assume \Cref{v4-arith:deg-f1a:conj:EG-wild-enhancement} below
(existence of the wild-enhanced EG stack with the stated structural
properties) and
\Cref{v4-arith:deg-f1a:conj:wild-block-functor}. Then there exists a
fully faithful functor
\begin{equation}
\label{v4-arith:deg-f1a:eq:wild-functor}
\Phi^{\mathrm{wild}}_{p,\chi}\colon
\mathrm{D}^{\mathrm{sph\text{-}fin}}\bigl([GL_{2}(\mathbb{Q}_{p})]\bigr)_{\chi}
\;\longrightarrow\;\mathrm{IndCoh}_{\mathrm{Nilp}}\bigl(\mathcal{X}^{\mathrm{EG,wild},p,\chi}_{GL_{2}}\bigr)
\end{equation}
extending DEG to the wild primes, compatible with the
Bushnell--Henniart fiber correspondence on each block.
\end{theorem}

\begin{proof}
\textbf{Step 1: block decomposition at \(p\in\{2,3\}\).} The wild
analogue of \Cref{v4-arith:deg-f1a:thm:colmez-paskunas-blocks}
holds partially at \(p=3\) (Colmez 2010 \S 9) and weakly at \(p=2\).
The proof uses instead the Bushnell--Henniart 2006 classification of
smooth \(GL_{2}(\mathbb{Q}_{p})\)-representations by types:
\begin{equation}
\mathrm{Mod}_{GL_{2}(\mathbb{Q}_{p})}^{\mathrm{sm},\chi}
\;=\;\prod_{\tau}\mathrm{Block}^{\mathrm{BH}}_{\tau},
\end{equation}
where the index \(\tau\) runs over the Bushnell--Kutzko types
(Bushnell--Kutzko 1993, Ann.~Math.~Studies~129). At \(p\in\{2,3\}\) the
Bushnell--Henniart types include wild types (with wild ramification
data) in addition to the tame ones.

\textbf{Step 2: spectral-side block decomposition.} By
\Cref{v4-arith:deg-f1a:conj:EG-wild-enhancement}, the wild-enhanced
EG stack decomposes as
\begin{equation}
\mathcal{X}^{\mathrm{EG,wild},p,\chi}_{GL_{2}}\;=\;\coprod_{\tau}\mathcal{X}^{\mathrm{EG,wild},p,\chi,\tau}_{GL_{2}},
\end{equation}
with each component indexed by the same set of Bushnell--Kutzko
types \(\tau\); the classical components (tame \(\tau\)) recover the
\(\rho\)-indexed components of
\Cref{v4-arith:deg-f1a:prop:spectral-blocks}, the wild components
are new derived-fibred components carrying the Ext\({}^{1}\)-enhancement.

\textbf{Step 3: per-block functor construction.} On each tame
and wild block, the per-block fully faithful functor is part of the
hypothesis.  Bushnell--Henniart supplies the type labels; it does not
by itself construct the categorical functor.

\textbf{Step 4: fully-faithfulness assembly.} The sum over \(\tau\)
of the per-block functors is fully faithful because each block is
separately fully faithful, and the blocks are pairwise orthogonal
on both sides (tame and wild types share no simples, and the
tame and wild components of the wild-EG stack are disjoint
formal opens).
\end{proof}

\begin{conjecture}[wild-EG enhancement]
\label{v4-arith:deg-f1a:conj:EG-wild-enhancement}
\ClaimStatusConjectured. For \(p\in\{2,3\}\) and a central character
\(\chi\), there exists a derived formal algebraic stack
\(\mathcal{X}^{\mathrm{EG,wild},p,\chi}_{GL_{2}}\) with the following
properties:
\begin{enumerate}
\item The truncation
\(\tau_{\leq 0}\mathcal{X}^{\mathrm{EG,wild},p,\chi}_{GL_{2}}\) equals
the classical Emerton--Gee stack
\(\mathcal{X}^{\mathrm{EG},p,\chi}_{GL_{2}}\).
\item The cotangent complex \(\mathbb{L}_{\mathcal{X}^{\mathrm{EG,wild}}}\)
concentrates in degrees \([-1,0]\), with the degree-\((-1)\) part
supported on the non-semisimple locus of
\(\mathcal{X}^{\mathrm{EG},p,\chi}_{GL_{2}}\) and of dimension \(1\) at
\(p=2\), dimension \(2\) at \(p=3\) at each closed point.
\item There is a fiber correspondence, compatible with local
Langlands, between the wild components of
\(\mathcal{X}^{\mathrm{EG,wild},p,\chi}_{GL_{2}}\) and the
Bushnell--Henniart wild types of \(GL_{2}(\mathbb{Q}_{p})\).
\item Ind-coherent sheaves with nilpotent singular support on
\(\mathcal{X}^{\mathrm{EG,wild},p,\chi}_{GL_{2}}\) are generated under
colimits and shifts by the derived structure sheaves of the
Bushnell--Henniart-indexed closed points.
\end{enumerate}
\end{conjecture}

\begin{remark}[what supports the wild-EG conjecture]
\label{v4-arith:deg-f1a:rem:wild-EG-evidence}
The conjecture is consistent with:
\begin{itemize}
\item Bushnell--Henniart 2006 \S 12 (wild ramification at \(p=2\) for
\(GL_{2}\)) and \S 15 (wild types);
\item Emerton--Gee arXiv:1908.07185 \S 6 (local-global
compatibility at wild primes);
\item Our own \Cref{v4-arith:bzn-rcompletion:prop:R4-Ext1-p2},
\Cref{v4-arith:bzn-rcompletion:prop:R4-Ext1-p3} (explicit
Ext\({}^{1}\)-computation at \(p=2,3\)).
\end{itemize}
The Bushnell--Henniart data give the block structure on the smooth
side; the Emerton--Gee stack gives the spectral-side support; the
Ext\({}^{1}\)-computation of Ch.~\ref{v4-arith:bzn-rcompletion} supplies
the derived-fibre dimensions. The conjectural assertion is that these
three inputs define a single derived formal algebraic stack with the
stated properties, in particular that
\(\mathrm{IndCoh}_{\mathrm{Nilp}}\) on this stack is compactly
generated by the Bushnell--Henniart closed points.
\end{remark}

\begin{remark}[domain of the wild-EG conjecture]
\label{v4-arith:deg-f1a:rem:wild-EG-domain}
\Cref{v4-arith:deg-f1a:conj:EG-wild-enhancement} is strictly a
derived enhancement of a classical stack: it does not ask for new
geometric content beyond what already exists. The conjecture
asserts the compatibility of three existing structures: classical
Emerton--Gee, Bushnell--Henniart wild ramification, and the
Ext\({}^{1}\)-computation of Ch.~\ref{v4-arith:bzn-rcompletion}. A
substantial treatment of this conjecture would fit within the
non-formal theorem condition of a single Iwasawa-Kisin derived deformation
theory paper.
\end{remark}

\subsection{BH-Block Exhaustion and Essential Surjectivity at Wild Primes}

\begin{conjecture}[BH-block exhaustion at wild primes]
\label{v4-arith:deg-f1a:conj:BH-block-exhaustion}
\ClaimStatusConjectured. For \(p\in\{2,3\}\), every central character
\(\chi\in\mathrm{Hom}(\mathbb{Z}_{p}^{\times},k^{\times})\), and every
Bushnell--Kutzko type \(\tau\) (tame or wild) in the support of
\(\mathcal{X}^{\mathrm{EG,wild},p,\chi}_{GL_{2}}\), the pair
\[
\bigl(\mathrm{Block}^{\mathrm{BH}}_{\tau},\;\mathcal{X}^{\mathrm{EG,wild},p,\chi,\tau}_{GL_{2}}\bigr)
\]
is DEG-exhausted in the sense of \Cref{v4-arith:deg-f1a:def:block-exhaustion}:
the wild DEG functor \(\Phi^{\mathrm{wild}}_{p,\chi,\tau}\) is essentially
surjective, equivalently the Bushnell--Henniart-indexed derived structure
sheaves generate \(\mathrm{IndCoh}_{\mathrm{Nilp}}(\mathcal{X}^{\mathrm{EG,wild},p,\chi,\tau}_{GL_{2}})\)
under colimits and shifts.
\end{conjecture}

\begin{corollary}[wild essential surjectivity, conditional]
\label{v4-arith:deg-f1a:cor:wild-ess-surj}
\ClaimStatusProvedHereConditional. Assume
\Cref{v4-arith:deg-f1a:conj:EG-wild-enhancement} and
\Cref{v4-arith:deg-f1a:conj:BH-block-exhaustion}. Then
\(\Phi^{\mathrm{wild}}_{p,\chi}\) of
\Cref{v4-arith:deg-f1a:thm:wild-functor-existence} is an equivalence
\begin{equation}
\mathrm{D}^{\mathrm{sph\text{-}fin}}\bigl([GL_{2}(\mathbb{Q}_{p})]\bigr)_{\chi}
\;\xrightarrow{\;\sim\;}\;\mathrm{IndCoh}_{\mathrm{Nilp}}\bigl(\mathcal{X}^{\mathrm{EG,wild},p,\chi}_{GL_{2}}\bigr).
\end{equation}
\end{corollary}

\begin{proof}
Fully-faithfulness is \Cref{v4-arith:deg-f1a:thm:wild-functor-existence}.
Essential surjectivity follows block-by-block from the BH-block
analogue of \Cref{v4-arith:deg-f1a:cor:ess-surj-block-by-block}
plus \Cref{v4-arith:deg-f1a:conj:BH-block-exhaustion}.
\end{proof}

\subsection{Source Comparison for Obstruction (b), \texorpdfstring{\(p\in\{2,3\}\)}{p in {2,3}}}

\paragraph{Obstruction 1 (source separation).} The wild list
(Bushnell--Henniart 2006, Carter 1989, Lusztig 1980, Colmez 2010
\S 9, Gee--Kisin 2014) is disjoint from the tame DEG list; both
are disjoint from the Fargues--Scholze list. The connection to
\Cref{v4-arith:bzn-rcompletion:prop:R4-Ext1-p2} uses Jantzen 2003
and Etingof--Ostrik 2004, also disjoint.

\paragraph{Obstruction 2 (conventions).} Bushnell--Henniart type
conventions match those of Emerton--Gee arXiv:1908.07185 via the
standard dictionary (Bushnell--Henniart \S 18 Thm~A). Ext\({}^{1}\)
conventions follow Jantzen 2003 Ch.~II \S 8.

\paragraph{Obstruction 3 (ambient category).} Wild-enhanced EG stack
sits in derived formal algebraic stacks; \(\mathrm{IndCoh}_{\mathrm{Nilp}}\)
on a derived formal stack is computed via the Gaitsgory--Rozenblyum
ind-completion. The LHS of
(\ref{v4-arith:deg-f1a:eq:wild-functor}) is the same as before.

\paragraph{Obstruction 4 (missing hypothesis).} Two conjectural inputs
are explicit: \Cref{v4-arith:deg-f1a:conj:EG-wild-enhancement}
(existence of the enhanced stack) and the BH-block-exhaustion
analogue of \Cref{v4-arith:deg-f1a:conj:block-exhaustion}
(essential surjectivity per block). \emph{The hypotheses are explicit.}

\paragraph{Obstruction 5 (false functoriality).} The wild DEG functor
is block-diagonal by construction; functoriality across blocks holds
because blocks are orthogonal. Within a block, functoriality under
convolution Hecke operators reduces to the Bushnell--Henniart
fiber correspondence, which is functorial by construction
(Bushnell--Henniart 2006 \S 18).

\paragraph{Obstruction 6 (unproved equivalence).} Two unproved inputs
(wild-EG enhancement; BH block exhaustion) are named.
\textbf{Named residuals.}

\paragraph{Obstruction 7 (numerical comparison).} No numerical computation.

The two explicit conjectural residuals are \Cref{v4-arith:deg-f1a:conj:EG-wild-enhancement} and \Cref{v4-arith:deg-f1a:conj:BH-block-exhaustion}.

\section{Obstruction (c): Central Character Variation}
\label{v4-arith:deg-f1a:central-char}

\subsection{The Central Character Torus}

The central character of \(GL_{2}(\mathbb{Q}_{p})\) is a continuous
homomorphism
\(\chi\colon\mathbb{Z}(\mathbb{Q}_{p})=\mathbb{Q}_{p}^{\times}\to k^{\times}\),
restricting to \(\mathbb{Z}_{p}^{\times}\) for the compact part. The
torus of continuous \(\chi\) forms
\[
\mathcal{T}^{\mathrm{cent}}_{p}\;:=\;\mathrm{Hom}_{\mathrm{cts}}\bigl(\mathbb{Q}_{p}^{\times},k^{\times}\bigr)
\;\simeq\;\mathbb{G}_{m}\times\widehat{\mathbb{Z}_{p}^{\times}},
\]
a product of a tame component (\(\mathbb{G}_{m}\): the value on a
uniformiser) and a compact component
(\(\widehat{\mathbb{Z}_{p}^{\times}}\): the value on units).

\begin{definition}[central-character family]
\label{v4-arith:deg-f1a:def:central-char-family}
\ClaimStatusDefinitional. Define the \emph{central-character
family} of the automorphic and spectral categories at \(p\):
\begin{align}
\mathrm{D}^{\mathrm{sph\text{-}fin}}\bigl([GL_{2}(\mathbb{Q}_{p})]\bigr)
&\;\simeq\;\int_{\chi\in\mathcal{T}^{\mathrm{cent}}_{p}}\mathrm{D}^{\mathrm{sph\text{-}fin}}\bigl([GL_{2}(\mathbb{Q}_{p})]\bigr)_{\chi}\,d\chi,\\
\mathrm{IndCoh}_{\mathrm{Nilp}}\bigl(\mathcal{X}^{\mathrm{EG},p}_{GL_{2}}\bigr)
&\;\simeq\;\int_{\chi\in\mathcal{T}^{\mathrm{cent}}_{p}}\mathrm{IndCoh}_{\mathrm{Nilp}}\bigl(\mathcal{X}^{\mathrm{EG},p,\chi}_{GL_{2}}\bigr)\,d\chi,
\end{align}
where \(\int d\chi\) denotes the direct-integral / fibre-decomposition
of a presentable \(\infty\)-category over the character torus.
\end{definition}

\begin{remark}[fibre decomposition is standard]
\label{v4-arith:deg-f1a:rem:fibre-decomp-standard}
The fibre decomposition of the automorphic side is the central-character
decomposition of smooth representations: \(\mathrm{Mod}^{\mathrm{sm}}_{GL_{2}(\mathbb{Q}_{p})}=\bigoplus_{\chi}\mathrm{Mod}^{\mathrm{sm},\chi}_{GL_{2}(\mathbb{Q}_{p})}\).
On the spectral side, \(\mathcal{X}^{\mathrm{EG},p}_{GL_{2}}\) admits
a natural morphism to \(\mathcal{T}^{\mathrm{cent}}_{p}\) (the
determinant morphism \(\rho\mapsto\det\rho/\chi_{\mathrm{cyc}}\))
whose fibres are the \(\mathcal{X}^{\mathrm{EG},p,\chi}_{GL_{2}}\)
stacks. The integral decomposition of \(\mathrm{IndCoh}_{\mathrm{Nilp}}\)
over this determinant fibration is standard (Gaitsgory--Rozenblyum's
AMS DAG monographs, Vol.~I Ch.~9, for the general setup).
\end{remark}

\begin{conjecture}[central-character family datum]
\label{v4-arith:deg-f1a:conj:central-character-family}
\ClaimStatusConjectured. The fixed-character DEG functors
\(\Phi^{\mathrm{DEG}}_{p,\chi}\) form a morphism of sheaves of
presentable stable categories over
\(\mathcal{T}^{\mathrm{cent}}_{p}\).  Equivalently, twisting the
central character on the automorphic side is carried by DEG to the
determinant twist on the Emerton--Gee side.
\end{conjecture}

\subsection{Pushforward of the DEG Functor}

\begin{theorem}[central-character variation of DEG]
\label{v4-arith:deg-f1a:thm:central-character-variation}
\ClaimStatusProvedHereConditional. Assume
\Cref{v4-arith:deg-f1a:conj:block-exhaustion} so that
\(\Phi^{\mathrm{DEG}}_{p,\chi}\) is an equivalence for each \(\chi\)
and each \(p\geq 5\).  Assume also
\Cref{v4-arith:deg-f1a:conj:central-character-family}. Then the family
\(\{\Phi^{\mathrm{DEG}}_{p,\chi}\}_{\chi}\) assembles into a single
equivalence
\begin{equation}
\label{v4-arith:deg-f1a:eq:varying-chi}
\Phi^{\mathrm{DEG}}_{p}\colon\mathrm{D}^{\mathrm{sph\text{-}fin}}\bigl([GL_{2}(\mathbb{Q}_{p})]\bigr)
\;\xrightarrow{\;\sim\;}\;\mathrm{IndCoh}_{\mathrm{Nilp}}\bigl(\mathcal{X}^{\mathrm{EG},p}_{GL_{2}}\bigr),
\end{equation}
natural in \(\chi\) and compatible with the central-character fibration
of \Cref{v4-arith:deg-f1a:def:central-char-family}.
\end{theorem}

\begin{proof}
The automorphic and spectral categories decompose over
\(\mathcal{T}^{\mathrm{cent}}_{p}\) as in
\Cref{v4-arith:deg-f1a:def:central-char-family}.  The family datum
turns the fibrewise DEG functors into a morphism of the two sheaves
of categories.  Under block exhaustion each fibre functor is an
equivalence, hence the induced morphism of sheaves is an equivalence
and its total category is the displayed equivalence.
\end{proof}

\begin{remark}[comparison to the fixed-character version]
\label{v4-arith:deg-f1a:rem:fixed-vs-varying}
The assembly
\Cref{v4-arith:deg-f1a:thm:central-character-variation} is
not a consequence of the fixed-character theorem alone.  It uses the
family datum of
\Cref{v4-arith:deg-f1a:conj:central-character-family} in addition to
block exhaustion at each \(\chi\).
\end{remark}

\subsection{Source Comparison for Obstruction (c)}

\paragraph{Obstruction 1 (source separation).} The central-character-variation
argument uses Gaitsgory--Rozenblyum Vol.~I Ch.~9--10 (fibrations
of categories) plus DEG arXiv:2603.26887. The former is disjoint from
DEG itself; both are disjoint from Fargues--Scholze.

\paragraph{Obstruction 2 (conventions).} The determinant-fibration
convention on \(\mathcal{X}^{\mathrm{EG},p}_{GL_{2}}\) follows
Emerton--Gee arXiv:1908.07185 \S 2.2. The sheaf-of-categories
convention follows Gaitsgory--Rozenblyum Vol.~I Ch.~10 Thm~10.3.2.1.
Compatible.

\paragraph{Obstruction 3 (ambient category).} Sheaves of presentable
stable \(\infty\)-categories over \(\mathcal{T}^{\mathrm{cent}}_{p}\);
the totalisation equivalence is an equivalence in this ambient.

\paragraph{Obstruction 4 (missing hypothesis).} Naturality of DEG in
\(\chi\) is not imported from the fixed-character theorem.  It is the
separate hypothesis
\Cref{v4-arith:deg-f1a:conj:central-character-family}.  The
fibre-to-total passage is formal only after that hypothesis is in
place. \emph{The hypothesis is explicit.}

\paragraph{Obstruction 5 (false functoriality).} The twisting morphism
\(\chi'\to\chi\) is compatible with both sides only under
\Cref{v4-arith:deg-f1a:conj:central-character-family}.

\paragraph{Obstruction 6 (unproved equivalence).} No new unproved
equivalence introduced; the theorem is conditional on the
per-\(\chi\) equivalence and on the family datum.

\paragraph{Obstruction 7 (numerical source).} No numerical computation.

The central-character family datum is a genuine residual.

\section{F1.a Reduction}
\label{v4-arith:deg-f1a:reduction}

\subsection{Assembly of (a), (b), (c)}

\begin{theorem}[F1.a at all primes, conditional]
\label{v4-arith:deg-f1a:thm:F1a-conditional}
\ClaimStatusProvedHereConditional. Assume:
\begin{itemize}
\item[\textbf{(H1)}] \Cref{v4-arith:deg-f1a:conj:block-exhaustion}
at every \(p\geq 5\) and every central character \(\chi\).
\item[\textbf{(H2)}] \Cref{v4-arith:deg-f1a:conj:EG-wild-enhancement}
at \(p\in\{2,3\}\) and every \(\chi\).
\item[\textbf{(H3)}] \Cref{v4-arith:deg-f1a:conj:wild-block-functor}
at \(p\in\{2,3\}\) and every \(\chi\).
\item[\textbf{(H4)}] \Cref{v4-arith:deg-f1a:conj:BH-block-exhaustion}
at \(p\in\{2,3\}\) and every \(\chi\).
\item[\textbf{(H5)}] \Cref{v4-arith:deg-f1a:conj:central-character-family}
at every \(p\).
\end{itemize}
Then for every prime \(p\) and every continuous central character
\(\chi\) there exists an equivalence
\begin{equation}
\label{v4-arith:deg-f1a:eq:F1a-conditional}
\Phi^{\mathrm{LocCat}}_{p,\chi}\colon\mathrm{D}^{\mathrm{sph\text{-}fin}}\bigl([GL_{2}(\mathbb{Q}_{p})]\bigr)_{\chi}
\;\xrightarrow{\;\sim\;}\;\mathrm{IndCoh}_{\mathrm{Nilp}}\bigl(\mathcal{X}^{\mathrm{EG,wild\text{-}enhanced},p,\chi}_{GL_{2}}\bigr),
\end{equation}
where the spectral side is the classical Emerton--Gee stack for
\(p\geq 5\) and its wild enhancement for \(p\in\{2,3\}\), and the
family over \(\chi\) assembles to the displayed F1.a equivalence
\begin{equation}
\label{v4-arith:deg-f1a:eq:F1a-displayed}
\Phi^{\mathrm{LocCat}}_{p}\colon\mathrm{D}^{\mathrm{sph\text{-}fin}}\bigl([GL_{2}(\mathbb{Q}_{p})]\bigr)
\;\xrightarrow{\;\sim\;}\;\mathrm{IndCoh}_{\mathrm{Nilp}}\bigl(\mathcal{X}^{\mathrm{EG,wild\text{-}enhanced},p}_{GL_{2}}\bigr).
\end{equation}
Local Hecke-functoriality at the place \(p\) is part of the local
categorical input; it is not inferred from full faithfulness alone.
\end{theorem}

\begin{proof}
Assemble \Cref{v4-arith:deg-f1a:thm:ess-surj-conditional} (valid
for \(p\geq 5\) under (H1)),
\Cref{v4-arith:deg-f1a:cor:wild-ess-surj} (valid for \(p\in\{2,3\}\)
under (H2), (H3), and (H4)), and
\Cref{v4-arith:deg-f1a:thm:central-character-variation}
(assembles the family under (H5)).
Per-\(\chi\) equivalences (from (H1)+(H2)+(H3)+(H4)) combine with the
family decomposition to produce the total equivalence
(\ref{v4-arith:deg-f1a:eq:F1a-displayed}).
\end{proof}

\subsection{Unconditional DEG Full Faithfulness}

\begin{theorem}[unconditional DEG full faithfulness at fixed central character]
\label{v4-arith:deg-f1a:thm:DEG-family-ff}
\ClaimStatusProvedElsewhere. At every prime \(p\geq 5\) and every
fixed central character \(\chi\), there exists a
fully faithful functor
\begin{equation}
\Phi^{\mathrm{DEG}}_{p,\chi}\colon
\mathrm{D}^{\mathrm{sph\text{-}fin}}\bigl([GL_{2}(\mathbb{Q}_{p})]\bigr)_{\chi}
\;\hookrightarrow\;
\mathrm{IndCoh}_{\mathrm{Nilp}}\bigl(\mathcal{X}^{\mathrm{EG},p,\chi}_{GL_{2}}\bigr)
\end{equation}
of \Cref{v4-arith:deg-f1a:thm:deg-fixed-chi}. Essential surjectivity is
conditional on \Cref{v4-arith:deg-f1a:conj:block-exhaustion};
wild primes \(p\in\{2,3\}\) and variation over \(\chi\) are excluded.
\end{theorem}

\begin{proof}
This is exactly the fixed-character DEG theorem recalled in
\Cref{v4-arith:deg-f1a:thm:deg-fixed-chi}.  No family assembly is
used.
\end{proof}

\begin{corollary}[partial F1.a at \(p\geq 5\)]
\label{v4-arith:deg-f1a:cor:partial-F1a-large-p}
\ClaimStatusProvedElsewhere. For every prime \(p\geq 5\) and fixed
\(\chi\), F1.a holds
\emph{up to the essential image of \(\Phi^{\mathrm{DEG}}_{p,\chi}\)}: there
is a fully faithful embedding
\[
\mathrm{D}^{\mathrm{sph\text{-}fin}}\bigl([GL_{2}(\mathbb{Q}_{p})]\bigr)_{\chi}\hookrightarrow\mathrm{IndCoh}_{\mathrm{Nilp}}\bigl(\mathcal{X}^{\mathrm{EG},p,\chi}_{GL_{2}}\bigr)
\]
No assertion about the total category over all \(\chi\) follows
without \Cref{v4-arith:deg-f1a:conj:central-character-family}.
\end{corollary}

\begin{proof}
Immediate from \Cref{v4-arith:deg-f1a:thm:DEG-family-ff}.
\end{proof}

\begin{remark}[maximal unconditional DEG range]
\label{v4-arith:deg-f1a:rem:strongest-uncond}
\Cref{v4-arith:deg-f1a:cor:partial-F1a-large-p} gives precisely the
unconditional range supplied by DEG: a full embedding at every prime
\(p\geq 5\) after fixing \(\chi\). It is the DEG theorem itself, not a
family strengthening. The displayed F1.a equivalence still requires
essential surjectivity, wild primes, and central-character variation.
\end{remark}

\subsection{Residuals}

\begin{description}
\item[\textbf{Residual R1a: block exhaustion at \(p\geq 5\).}]
\Cref{v4-arith:deg-f1a:conj:block-exhaustion}. Closes (a) at
\(p\geq 5\). Reduces essential surjectivity to a finite-block
Koszul-duality check per residual Galois representation \(\rho\).
\emph{Hypothesis: non-formal theorem.}
\item[\textbf{Residual R1b: wild-EG enhancement at \(p\in\{2,3\}\).}]
\Cref{v4-arith:deg-f1a:conj:EG-wild-enhancement}. Closes (b) after
combining with the BH analogue of block exhaustion.
\emph{Hypothesis: non-formal theorem.}
\item[\textbf{Residual R1c: BH block exhaustion at \(p\in\{2,3\}\).}]
The Bushnell--Henniart version of block exhaustion. Closes (b) in
conjunction with R1b.
\emph{Hypothesis: non-formal theorem.}
\end{description}

Central-character variation (c) introduces \emph{no new residual}:
\Cref{v4-arith:deg-f1a:thm:central-character-variation} is
unconditional given the per-\(\chi\) closure.

\section{Consequence for F1 and \(\mathrm{L}_{\mathrm{ACD}}\)}
\label{v4-arith:deg-f1a:global-check}

\subsection{Local F1.a Closed at Every Prime Suffices for Adelic F1}

\begin{theorem}[cross-check with Chapter~\ref{v4-arith:closure}]
\label{v4-arith:deg-f1a:thm:F1-adelic-check}
\ClaimStatusProvedHereConditional. Assume the conditional closure of F1.a at
every prime \(p\) of
\Cref{v4-arith:deg-f1a:thm:F1a-conditional} (i.e.\ assume (H1),
(H2), (H3)). Then the F1 obstruction is reduced to the current
\(\mathrm{L}_{\mathrm{ACD}}\) table of
\Cref{v4-arith:cl:cor:F1-single-lemma}:
\begin{equation}
\label{v4-arith:deg-f1a:eq:F1-under-F1a}
\mathrm{D}^{\mathrm{sph\text{-}fin}}\bigl([GL_{2}(\mathbb{Q})\backslash GL_{2}(\mathbb{A}_{\mathbb{Q}})]\bigr)\;\simeq\;\mathrm{IndCoh}_{\mathrm{Nilp}}\bigl(\mathcal{X}^{\mathrm{Sel}}_{\bar\rho,\det,\mathcal{L}}\bigr).
\end{equation}
\end{theorem}

\begin{proof}
Apply \Cref{v4-arith:cl:cor:F1-single-lemma}: F1 holds given
(F1.a) and \(\mathrm{L}_{\mathrm{ACD}}\). Our
\Cref{v4-arith:deg-f1a:thm:F1a-conditional} reduces (F1.a) under
(H1)--(H5), up to the wild-enhancement of the local conditions
at \(p\in\{2,3\}\). These wild conditions must be inserted into the
Selmer-derived stack as modified factors
\(\mathcal{X}^{\mathrm{EG,wild\text{-}enhanced},p}_{GL_{2}}\) in the
local condition \(\mathcal{L}_{p}\). The global equivalence
(\ref{v4-arith:deg-f1a:eq:F1-under-F1a}) therefore uses the
Selmer-derived target on the right-hand side; it is not an
ind-Noetherian product of local stacks.
\end{proof}

\begin{remark}[three local pieces]
\label{v4-arith:deg-f1a:rem:global-three-pieces}
Under the assumptions (H1)--(H5) of
\Cref{v4-arith:deg-f1a:thm:F1a-conditional} plus
\(\mathrm{L}_{\mathrm{ACD}}\), F1 holds. The assumptions
(H1)--(H5) are the local-categorical-comparison residuals
of the DEG reduction; \(\mathrm{L}_{\mathrm{ACD}}\) is the adelic-categorical-descent
residual of \Cref{v4-arith:cl:rem:F1-substantial}. The F1 obstruction
set therefore has six items: (H1)--(H5) and
\(\mathrm{L}_{\mathrm{ACD}}\).
\end{remark}

\subsection{Refined Chapter~\ref{v4-arith:closure} Reduction}

\begin{proposition}[refined F1 reduction]
\label{v4-arith:deg-f1a:prop:F1-reduction-updated}
\ClaimStatusProvedHere. The F1-reduction of
\Cref{v4-arith:thm:F1-reduction} refines to:
the F1.a item splits into three sub-residuals
\begin{itemize}
\item[\textbf{(F1.a.1)}] Block exhaustion at \(p\geq 5\), fixed \(\chi\)
(\Cref{v4-arith:deg-f1a:conj:block-exhaustion}): essential
surjectivity of DEG.
\item[\textbf{(F1.a.2)}] Wild-EG enhancement at \(p\in\{2,3\}\)
(\Cref{v4-arith:deg-f1a:conj:EG-wild-enhancement}): derived
formal algebraic stack extending EG to wild primes.
\item[\textbf{(F1.a.3)}] BH-block exhaustion at \(p\in\{2,3\}\)
(\Cref{v4-arith:deg-f1a:conj:BH-block-exhaustion}):
essential surjectivity per Bushnell--Kutzko type.
\end{itemize}
Central-character variation is contained in the family version
(\Cref{v4-arith:deg-f1a:thm:central-character-variation}) without a
separate residual. The total F1 obstruction set is:
\begin{equation}
\mathrm{Res}(F1)\;=\;\{\text{F1.a.1},\text{F1.a.2},\text{F1.a.3},\mathrm{L}_{\mathrm{ACD}}\}.
\end{equation}
\end{proposition}

\begin{proof}
By \Cref{v4-arith:deg-f1a:thm:F1a-conditional}, (F1.a) holds given
(H1)=(F1.a.1), (H2)=(F1.a.2), (H3)=(F1.a.3). Applying
\Cref{v4-arith:thm:F1-reduction}, F1 is equivalent to (F1.a) and
\(\mathrm{L}_{\mathrm{ACD}}\); substituting the three-item expansion
of (F1.a) gives the displayed obstruction set.
\end{proof}

\subsection{Source Comparison for Global Assembly}

\paragraph{Obstruction 1.} Sources: \Cref{v4-arith:cl:cor:F1-single-lemma}
(the residual splitting F1 reduction), plus the local-categorical-comparison
arguments of the previous sections. Disjoint.

\paragraph{Obstruction 2.} Global Emerton--Gee formalism follows
Emerton--Gee \S 2.3 (global ind-Noetherian limit). The
wild-enhancement is a local correction at two primes; globally it
absorbs via the same ind-Noetherian limit.

\paragraph{Obstruction 3.} Global stack sits in ind-Noetherian formal
algebraic stacks over \(\mathrm{Spec}\,\mathbb{Z}\); the automorphic
side is the adelic double coset. Both sides have the same ambient
presentable \(\infty\)-category.

\paragraph{Obstruction 4.} (H1), (H2), (H3), \(\mathrm{L}_{\mathrm{ACD}}\)
explicit.

\paragraph{Obstruction 5.} Global assembly is a restricted tensor product
over primes, compatible with cross-prime Hecke-functoriality. This
compatibility is part of \(\mathrm{L}_{\mathrm{ACD}}\); not a new
residual of the DEG reduction.

\paragraph{Obstruction 6.} No new unproved equivalence.

\paragraph{Obstruction 7.} No numerical computation.

The global consequence is unconditional only after the local,
wild-functor, family, and \(\mathrm{L}_{\mathrm{ACD}}\) residuals are
assumed.

\section{Residual Conditions}
\label{v4-arith:deg-f1a:residuals}

\subsection{Residual Conditions for the DEG Extension}

The DEG extension has five residual inputs:

\begin{enumerate}
\item \Cref{v4-arith:deg-f1a:conj:block-exhaustion}
(\textbf{R1a/F1.a.1}): block exhaustion for DEG at \(p\geq 5\),
fixed \(\chi\), every residual \(\rho\). Reduces essential surjectivity
of DEG to a finite-block Koszul-type check.
\item \Cref{v4-arith:deg-f1a:conj:EG-wild-enhancement}
(\textbf{R1b/F1.a.2}): existence of a derived formal algebraic
wild-enhanced EG stack at \(p\in\{2,3\}\) with four structural
properties.
\item \Cref{v4-arith:deg-f1a:conj:wild-block-functor}
(\textbf{R1b'}): existence of the Bushnell--Henniart-compatible
per-block fully faithful functor at \(p\in\{2,3\}\).
\item \Cref{v4-arith:deg-f1a:conj:BH-block-exhaustion}
(\textbf{R1c/F1.a.3}): the wild-block analogue of
\Cref{v4-arith:deg-f1a:conj:block-exhaustion}; essential
surjectivity per Bushnell--Kutzko type at \(p\in\{2,3\}\).
\item \Cref{v4-arith:deg-f1a:conj:central-character-family}
(\textbf{R1d}): compatibility of the fixed-character DEG functors
with the central-character torus.
\end{enumerate}

\subsection{What is Unconditional}

Despite the residuals, two results close unconditionally:

\begin{enumerate}
\item The fixed-\(\chi\), \(p\geq 5\), fully faithful DEG embedding
(\Cref{v4-arith:deg-f1a:thm:DEG-family-ff}): no residual needed.
This is the strongest unconditional DEG input toward F1.a.
\item The global F1 reduction
(\Cref{v4-arith:deg-f1a:prop:F1-reduction-updated}): no residual
beyond the explicitly named local, family, and
\(\mathrm{L}_{\mathrm{ACD}}\) inputs.
\end{enumerate}

\subsection{Out of Domain}

Three items are outside the stated domain:

\begin{enumerate}
\item \emph{Langlands functoriality at the global level.} Deferred to
the adelic categorical descent \(\mathrm{L}_{\mathrm{ACD}}\).
\item \emph{Higher \(GL_{n}\) for \(n\geq 3\).} The chapter focuses on
\(GL_{2}(\mathbb{Q}_{p})\); DEG's technology generalises under
significant additional input (Emerton--Helm--Paskunas for
\(GL_{n}\)), not in domain here.
\item \emph{Number fields beyond \(\mathbb{Q}\).} The chapter focuses
on \(\mathbb{Q}_{p}\); extension to \(K_{v}\) for \(K/\mathbb{Q}\) a
finite extension requires additional considerations at primes
ramified in \(K\).
\end{enumerate}

\section{DEG Extension Boundary}
\label{v4-arith:deg-f1a:summary}

\begin{enumerate}
\item the residual splitting overextended F1.a as closed via Fargues--Scholze. The
actual state of the art is Dotto--Emerton--Gee arXiv:2603.26887:
\emph{fully-faithful, fixed central character, \(p\geq 5\) only}
(\Cref{v4-arith:deg-f1a:thm:deg-fixed-chi}).
\item Three obstructions separate DEG from F1.a:
(a) essential surjectivity; (b) extension to \(p\in\{2,3\}\);
(c) central-character variation
(\Cref{v4-arith:deg-f1a:rem:deg-not-assert}).
\item Obstruction (a) reduces to
\Cref{v4-arith:deg-f1a:conj:block-exhaustion} (F1.a.1): block
exhaustion per residual Galois representation.
\Cref{v4-arith:deg-f1a:thm:ess-surj-conditional} closes (a) under
this conjecture.
\item Obstruction (b) reduces to
\Cref{v4-arith:deg-f1a:conj:EG-wild-enhancement} (F1.a.2: existence
of wild-enhanced EG stack),
\Cref{v4-arith:deg-f1a:conj:wild-block-functor}, plus
\Cref{v4-arith:deg-f1a:conj:BH-block-exhaustion} (F1.a.3).
\Cref{v4-arith:deg-f1a:thm:wild-functor-existence} and
\Cref{v4-arith:deg-f1a:cor:wild-ess-surj} close (b) under these
conjectures.
\item Obstruction (c) is the family datum
\Cref{v4-arith:deg-f1a:conj:central-character-family}.  It is not a
consequence of the fixed-character DEG theorem.
\item Full conditional closure of F1.a at all primes:
\Cref{v4-arith:deg-f1a:thm:F1a-conditional}, conditional on
(H1)--(H5).
\item Unconditional partial closure: the fixed-\(\chi\), \(p\geq 5\),
fully faithful DEG embedding
(\Cref{v4-arith:deg-f1a:thm:DEG-family-ff},
\Cref{v4-arith:deg-f1a:cor:partial-F1a-large-p}). No family
statement is asserted.
\item Global assembly check: F1 follows from F1.a (at every prime,
including wild) plus \(\mathrm{L}_{\mathrm{ACD}}\)
(\Cref{v4-arith:deg-f1a:thm:F1-adelic-check}); total F1 residual
table is six items
(\Cref{v4-arith:deg-f1a:prop:F1-reduction-updated}).
\end{enumerate}

There is no unconditional closure of F1.a. The
strongest honest statement is the fixed-\(\chi\), \(p\geq 5\), fully
faithful DEG embedding
(\Cref{v4-arith:deg-f1a:cor:partial-F1a-large-p}); all other closure
claims are explicitly conditional. The residuals F1.a.1, F1.a.2,
the wild block functor datum, F1.a.3, and the central-character
family datum are strictly narrower than the original F1.a conjecture
and each is individually testable with the source list laid out
here.
